% ============================================================
%  Carta de Resposta aos Avaliadores — Revista Análise Econômica
%  Documento NÃO identificado (sem autores)
% ============================================================
\documentclass[a4paper,12pt]{article}

\usepackage[utf8]{inputenc}
\usepackage[T1]{fontenc}
\usepackage[brazil]{babel}
\usepackage[
  a4paper,
  top=3cm, bottom=3cm,
  left=3cm, right=3cm
]{geometry}
\usepackage{setspace}
\setstretch{1.3}
\setlength{\parindent}{1.25cm}
\setlength{\parskip}{4pt}
\usepackage{enumitem}
\usepackage{booktabs}
\usepackage{xcolor}
\usepackage{hyperref}
\usepackage{microtype}

% Cores para os avaliadores
\definecolor{corA}{RGB}{0,70,140}
\definecolor{corB}{RGB}{140,40,0}

\newcommand{\avA}[1]{\textcolor{corA}{\textit{#1}}}
\newcommand{\avB}[1]{\textcolor{corB}{\textit{#1}}}
\newcommand{\resposta}[1]{\noindent\textbf{Resposta:} #1}

\begin{document}

\begin{center}
  {\large\bfseries Carta de Resposta aos Avaliadores}\\[0.4em]
  {\normalsize Revista Análise Econômica}\\[0.2em]
  {\small Manuscrito: \textit{Emprego Formal e Desastres Climáticos:
  Evidências Causais das Enchentes de 2024 no Rio Grande do Sul}}
\end{center}

\bigskip

\noindent Prezados Editores e Avaliadores,

\medskip

Agradecemos os pareceres recebidos. Os comentários foram cuidadosos e revelam leitura
atenta dos aspectos metodológicos e dos limites de inferência do manuscrito original.
Em resposta, reformulamos integralmente a estratégia empírica: abandonamos o modelo
C-ARIMA e adotamos o \textit{Synthetic Difference-in-Differences} (SDiD), complementado
por TWFE-DiD e Controle Sintético como triangulação, construindo um contrafactual
explícito a partir de 193 municípios gaúchos com $\leq 0{,}5\%$ da população atingida.
A amostra de municípios tratados foi ampliada de dois para seis. As demais alterações
são descritas a seguir; comentários específicos ao C-ARIMA tornaram-se inaplicáveis ao
novo desenho e são assim indicados.

\bigskip
\hrule
\bigskip

\section*{Respostas ao Avaliador A}

\subsection*{Comentário A.1 — Introdução}

\avA{``A introdução precisa ser totalmente reescrita citando fontes em informações
fortes, conectando melhor as afirmações descritas, inserindo melhor o artigo na
literatura que se pretende, citando melhor a metodologia e inserindo qual é a fonte
de dados utilizada (lacuna presente inclusive no resumo/abstract do artigo).''}

\bigskip

\resposta{A introdução foi integralmente reescrita. O novo texto abre com estatísticas
precisas do desastre (Defesa Civil Estadual; IPEA, 2024), apresenta a literatura de
custos diretos e indiretos de inundações (Botzen, Deschenes e Sanders, 2019; Berlemann
e Wenzel, 2018), motiva a escolha das admissões formais com base em Teixeira et al.\
(2024) e justifica a adoção do SDiD em detrimento de métodos alternativos. A fonte de
dados (CAGED/Novo CAGED) é mencionada na introdução e no resumo/abstract. A seção
``Dados'' (Seção~4) complementa com detalhes sobre a construção da série e os critérios
de seleção da amostra.}

\bigskip

\subsection*{Comentário A.2 — Revisão de literatura insuficiente}

\avA{``A seção de revisão de literatura precisa ser complementada com artigos que tratem
de desastres climáticos, do estado do Rio Grande do Sul, e também da metodologia
empregada.''}

\bigskip

\resposta{A revisão de literatura (Seção~2) foi reestruturada em quatro blocos: (i) a
evolução dos métodos quase-experimentais em dados de painel para avaliação de desastres
(TWFE-DiD, Controle Sintético e SDiD), com referências a Cavallo et al.\ (2013),
Deryugina (2017, 2018), Leiter, Oberhofer e Raschky (2009) e Clarke et al.\ (2023),
entre outros; (ii) impacto econômico de inundações, incluindo mercado de trabalho,
custos diretos e indiretos, e evidências para firmas europeias; (iii) evidências sobre
agricultura, pobreza e desigualdade; e (iv) desafios de identificação e fronteiras
metodológicas, com discussão de SUTVA, tendências paralelas e heterogeneidade de
efeitos. O trabalho de Teixeira et al.\ (2024) sobre as enchentes de 2024 no RS é
discutido em detalhe em todas as seções relevantes.}

\bigskip

\subsection*{Comentário A.3 — Seção metodológica confusa}

\avA{``Na seção metodológica, precisa introduzir melhor os modelos utilizados e
reescrever os parágrafos confusos e inconclusivos.''}

\bigskip

\resposta{A seção metodológica foi integralmente substituída (Seção~3). Os três
estimadores (TWFE-DiD, Controle Sintético e SDiD) são apresentados com suas
formulações matemáticas, premissas de identificação e procedimentos de inferência.
Uma tabela comparativa (Tabela~1) sintetiza as diferenças entre os estimadores,
facilitando a leitura. O SDiD é apresentado como estratégia principal, com
justificativa explícita baseada na divergência esperada de pré-tendências entre
municípios afetados e o pool de doadores.}

\bigskip

\subsection*{Comentário A.4 — Testes de diferença de amostras, quebra estrutural,
KPSS, ADF na tendência, lags de Roca Sales}

\avA{``Precisa apresentar testes t de diferença de amostras, além de testes de
quebra-estrutural para as séries. [...] justificar a escolha do teste ADF e apresentar
resultados de KPSS. [...] estimar os modelos também na tendência. [...] modelos com
lag entre 4 e 12 no município de Roca Sales.''}

\bigskip

\resposta{Estes comentários são específicos ao modelo C-ARIMA, que não integra o artigo
revisado. A mudança para o SDiD dispensa os testes de raiz unitária (ADF, KPSS) e de
quebra estrutural como pré-requisito de estimação, pois o SDiD não requer
estacionaridade da série — trabalha em diferenças ponderadas e elimina tendências
fixas via efeitos fixos de tempo. A preocupação de fundo — a dinâmica atípica de
2020--2021 na série de Eldorado do Sul — é tratada pela janela de estimação que inicia
em janeiro de 2021 (pós-pandemia), decisão que garante a homogeneidade metodológica da
série (passagem do Caged Antigo para o Novo CAGED também ocorreu em 2020) e é
justificada na Seção~4.1. A robustez desta escolha é avaliada pelo banco de pesos
temporais $\hat{\lambda}_t$ estimados pelo SDiD, que concentra peso nos períodos
pré-tratamento mais informativos e reduz a influência de janelas atípicas sem excluí-las
arbitrariamente.}

\bigskip

\subsection*{Comentário A.5 — Linguagem causal na interpretação}

\avA{``A linguagem da interpretação dos resultados está causal demais, sobretudo a parte
em que relaciona os resultados do artigo a resultados diretos de políticas públicas
realizadas nos municípios.''}

\bigskip

\resposta{Este comentário foi integralmente acolhido. No artigo revisado, a análise de
heterogeneidade e canais de dano (Seção~5.5) é explicitamente qualificada como
\textit{descritiva} e \textit{geradora de hipóteses}, com advertência metodológica
destacada em caixa de recuo. A menção a políticas públicas específicas como
mecanismos de recuperação foi suprimida da seção de resultados e migrou para a agenda
de pesquisa nas Considerações Finais, onde é posicionada como extensão causal que
requer desenho empírico próprio — exatamente a recomendação do Avaliador~A. A
linguagem das Considerações Finais foi alinhada com a seção de resultados.}

\bigskip

\subsection*{Comentário A.6 — Referências infladas}

\avA{``As referências estão infladas, com artigos que não são citados ao longo do
texto.''}

\bigskip

\resposta{A bibliografia foi inteiramente reconstruída a partir do novo texto. Todas as
referências constantes na lista final são citadas ao longo do manuscrito. O banco
bibliográfico foi depurado.}

\bigskip
\hrule
\bigskip

\section*{Respostas ao Avaliador B}

\subsection*{Comentários B.1.1--B.1.3 — Justificativa para o método e ausência de
grupo de controle externo}

\avB{``A justificativa para a adoção do C-ARIMA baseia-se na limitação do DiD em
capturar efeitos médios [...] não é uma justificativa suficiente para substituir métodos
com grupo de controle externo. [...] Recomenda-se incluir, ao menos como teste de
robustez, um exercício em DiD ou Controle Sintético.''}

\bigskip

\resposta{Acolhido como recomendação central desta revisão — ver abertura desta carta
para a descrição geral da mudança. Em complemento, a versão C-ARIMA é mencionada na
Seção~2.5 como trabalho anterior que motivou a reformulação, atendendo à sugestão de
fortalecer o posicionamento frente à literatura.}

\bigskip

\subsection*{Comentário B.2 — Desbalanceamento na seção de estratégia empírica}

\avB{``A exposição do modelo ARIMA clássico ocupa espaço desproporcional em relação à
extensão causal do modelo [...] Sugere-se reequilibrar o espaço dedicado a cada
componente.''}

\bigskip

\resposta{Com a adoção do SDiD, a seção metodológica passou a ser estruturada de forma
equilibrada: cada estimador (TWFE-DiD, Controle Sintético e SDiD) recebe uma
subseção dedicada com formulação formal, premissas de identificação e procedimento de
inferência. A ênfase concentra-se no SDiD e em sua dupla robustez — a contribuição
metodológica central —, com tratamento comparativo que justifica a adoção em detrimento
dos estimadores alternativos.}

\bigskip

\subsection*{Comentários B.3.1--B.3.4 — Pressupostos para inferência causal}

\avB{``A discussão é predominantemente declaratória: afirma-se que cada pressuposto é
atendido sem que o argumento seja desenvolvido de forma analítica. [...] tensão entre
irreversibilidade do choque e recuperação rápida [...] ausência de interferência temporal
precisa de discussão explícita.''}

\bigskip

\resposta{As premissas de identificação do SDiD (tendências paralelas ponderadas,
não antecipação, SUTVA) são discutidas analiticamente na Seção~2.4 e verificadas
empiricamente. A tensão entre ``irreversibilidade'' e ``recuperação'' presente no
C-ARIMA não se reproduz no SDiD, cujo estimador captura o ATT médio pós-evento sem
presupor a persistência do choque. A premissa de não-interferência temporal (ausência
de outros eventos simultâneos) é abordada via zona de exclusão \textit{donut}
(tratados com $>50\%$ atingidos; doadores com $\leq 0{,}5\%$), que mitiga a
contaminação do contrafactual. A discussão sobre SUTVA e transbordamentos (cadeia
produtiva, deslocamento de demanda) é feita explicitamente nas Seções~2.4.3 e 4.3. A
quantidade de períodos pré-tratamento disponíveis ($T_0 = 40$ meses) é informada na
Seção~4.1, superando as preocupações sobre séries curtas levantadas pelo Avaliador~B.}

\bigskip

\subsection*{Comentário B.4 — Critério de seleção dos municípios}

\avB{``O critério objetivo de seleção não é apresentado diretamente: não há tabela,
ranking ou threshold que permita ao leitor verificar por que esses dois municípios foram
escolhidos entre os mais de 400 afetados.''}

\bigskip

\resposta{O critério de seleção foi tornando explícito e verificável: são tratados todos
os municípios gaúchos com mais de 50\% da população atingida segundo o MUPRS/SPGG-RS
— resultando em seis municípios (Eldorado do Sul, Muçum, Roca Sales, Travesseiro,
Arambaré e Igrejinha), apresentados na Tabela~2 com código IBGE, população e
proporção atingida. O limiar de 50\% é justificado como critério de \textit{alta
severidade relativa} (disrupção de escala municipal), pré-especificado antes da
estimação, e sua escolha é discutida na Seção~4.2. Amostra de dois municípios tornou-se
amostra de seis, aumentando o poder da análise e sua generalização para o contexto
do desastre.}

\bigskip

\subsection*{Comentário B.5 — Variável dependente e teste com saldo líquido}

\avB{``Rodar o mesmo modelo utilizando o saldo líquido de empregos deveria produzir
resultados consistentes — o que, se confirmado, fortaleceria consideravelmente os
achados.''}

\bigskip

\resposta{Este teste foi implementado integralmente no Apêndice~A. A triangulação
SDiD/SC/TWFE-DiD é repetida com o saldo líquido mensal (admissões $-$ desligamentos)
como variável dependente sobre o mesmo painel e pool de doadores. Os resultados
confirmam que o ajuste ao desastre se deu predominantemente via retração das
contratações: nos municípios com efeito robusto (Muçum, Roca Sales, Travesseiro), os
ATTs sobre o saldo são negativos e significantes, com magnitudes próximas aos ATTs
sobre as admissões. A exceção documentada — Eldorado do Sul, cujo SDiD para o saldo
é positivo ($+1{,}83$; $p = 0{,}057$) — é explicada pelo congelamento simultâneo de
contratações \textit{e} demissões no período agudo, padrão consistente com a rigidez
da legislação trabalhista brasileira e discutido à luz de Teixeira et al.\ (2024) e
Krein (2012).}

\bigskip

\subsection*{Comentário B.6 — Tratamento do período pandêmico}

\avB{``O argumento para manter o período pandêmico na série de treinamento é razoável,
mas a defesa é incompleta. [...] Sugere-se que o autor inclua os resultados obtidos
com e sem o período pandêmico.''}

\bigskip

\resposta{A janela de estimação foi reposicionada para iniciar em janeiro de 2021 por
dois motivos: (i) garantir homogeneidade metodológica da série após a transição do
Caged Antigo para o Novo CAGED em janeiro de 2020; e (ii) excluir a fase aguda da
pandemia (2020), cujo padrão atípico comprometeria a validade dos pesos sintéticos
estimados sobre o período pré-tratamento. Essa escolha substitui a análise de
sensibilidade solicitada pelo Avaliador~B: em vez de testar com e sem pandemia, o
problema foi resolvido diretamente pela escolha do ponto de início da série. Os pesos
de tempo $\hat{\lambda}_t$ do SDiD garantem robustez adicional ao atribuir menor peso
a períodos historicamente distantes ou atípicos.}

\bigskip

\subsection*{Comentários B.7 — Qualidade do ajuste; ambiguidade no termo ``recalibrar''}

\avB{``Para Eldorado do Sul, o RMSE de aproximadamente 121 admissões representa um
erro da ordem de 60\% da média da série. [...] O autor deve contextualizar as métricas
em relação à escala da série. [...] O uso do termo `recalibrar' gera ambiguidade.''}

\bigskip

\resposta{O diagnóstico de ajuste foi o principal motivador da mudança de estimador. O
SDiD minimiza por construção o erro quadrático ponderado pré-tratamento, produzindo um
contrafactual que replica a trajetória observada antes do evento. O diagnóstico adotado
no artigo revisado é o RMSPE pré-tratamento (Seção~5.2): para Eldorado do Sul, 2,39
admissões por mil habitantes; para Muçum (RMSPE $= 7{,}07$), a inferência é qualificada
como parcial. O termo ``recalibrar'' foi suprimido: o SDiD não tem etapas de
treinamento/teste sequenciais.}

\bigskip

\subsection*{Comentário B.8 — Interpretação da dinâmica pós-choque e atribuição a
políticas}

\avB{``A atribuição da recuperação a políticas públicas específicas extrapola o que o
modelo efetivamente estima. Recomenda-se que essa seção seja reposicionada como
discussão de `possíveis mecanismos'.''}

\bigskip

\resposta{Acolhido — ver resposta ao Comentário A.5. Nas Considerações Finais, a
avaliação de políticas de recuperação específicas é posicionada como agenda de pesquisa
futura, reconhecendo que sua identificação causal requer desenho empírico próprio.}

\bigskip

\subsection*{Comentário B.9 — Teste placebo: opacidade metodológica (kNN)}

\avB{``A principal questão é a opacidade metodológica na aplicação do algoritmo kNN
para seleção dos municípios placebo. [...] O leitor não consegue avaliar se os
municípios selecionados são de fato comparáveis.''}

\bigskip

\resposta{O teste placebo do artigo revisado é o teste de permutação direta no espaço
(placebo in-space), procedimento com validade exata para qualquer tamanho amostral
(Abadie, Diamond e Hainmueller, 2010). O procedimento consiste em estimar o SDiD para
cada um dos 193 municípios do pool doador como se fosse o tratado e comparar o efeito
real com essa distribuição de efeitos placebo. Não há seleção subjetiva por kNN; o
conjunto de municípios placebo é idêntico ao pool doador da estimação principal. Os
$p$-valores são calculados como a proporção de placebos com efeito ao menos tão extremo
quanto o efeito real, procedimento padrão na literatura de controle sintético e SDiD.
A transparência do procedimento é completa: o pool de 193 municípios é listado na
Seção~4.3, e o $p$-valor de cada município tratado é reportado na Tabela~1.}

\bigskip

\subsection*{Comentários B.10 — Comparação com Teixeira et al.\ (2024); limitações
não declaradas}

\avB{``A comparação com Teixeira et al.\ (2024) é apresentada como validação mútua,
sem considerar que os estudos utilizam métodos com pressupostos distintos. [...] A
limitação metodológica central não é mencionada nas Considerações Finais.''}

\bigskip

\resposta{A comparação com Teixeira et al.\ (2024) foi reposicionada: o artigo
revisado discute as diferenças de método (DiD vs.\ SDiD), amostra e normalização,
apresentando a convergência de sinal como evidência consistente — não como validação
mútua. As limitações são declaradas explicitamente nas Considerações Finais (Seção~6):
(i) restrição ao emprego formal; (ii) heterogeneidade descritiva com $n = 6$; (iii)
janela encerrada possivelmente antes da recuperação completa. A preocupação sobre
fatores simultâneos ao choque é endereçada pela zona de exclusão \textit{donut} e pela
discussão de SUTVA nas Seções~2.4.3 e~4.3.}

\bigskip
\hrule
\bigskip

\section*{Síntese das Alterações Realizadas}

\begin{table}[h!]
\centering
\small
\begin{tabular}{p{5.5cm} p{4.5cm} p{4cm}}
\toprule
\textbf{Crítica original} & \textbf{Alteração realizada} & \textbf{Avaliador(es)} \\
\midrule
Ausência de grupo de controle externo &
Adoção do SDiD com pool de 193 municípios doadores &
B (1.1--1.3) \\[4pt]
Critério de seleção dos municípios não explícito &
Seis municípios com $>50\%$ atingidos; Tabela com critério e MUPRS &
B (4.1) \\[4pt]
Sem teste de robustez com saldo líquido &
Triangulação replicada com saldo líquido (Apêndice~A) &
B (5.1) \\[4pt]
Revisão de literatura insuficiente &
Revisão reestruturada em quatro blocos (métodos, impacto, agricultura, identificação) &
A \\[4pt]
Introdução fraca; sem dados/método/fonte &
Introdução reescrita com estatísticas precisas, método, fonte de dados (CAGED) &
A \\[4pt]
Metodologia confusa e desequilibrada &
Seção 3 com formulação formal dos três estimadores; Tabela comparativa &
A, B (2) \\[4pt]
Linguagem causal; atribuição a políticas &
Análise qualificada como descritiva; políticas como agenda futura &
A, B (8) \\[4pt]
Placebo com kNN opaco &
Permutação direta sobre os 193 doadores; p-valor exato e transparente &
B (9) \\[4pt]
RMSE não contextualizado &
RMSPE pré-tratamento contextualizado; Muçum qualificado como inferência parcial &
B (7) \\[4pt]
Referências infladas &
Bibliografia reconstruída; apenas citações efetivas no texto &
A \\[4pt]
Período pandêmico insuficientemente tratado &
Série inicia em jan./2021; justificativa na Seção~4.1 &
B (6) \\[4pt]
Limitação metodológica central não declarada &
Limitações declaradas explicitamente; SUTVA discutido nas Seções~2.4.3 e 4.3 &
B (10.2) \\
\bottomrule
\end{tabular}
\end{table}

\bigskip

Reiteramos nossa gratidão aos dois avaliadores pela qualidade e profundidade das
revisões. A combinação dos pareceres foi instrumental para a profunda reformulação que
o trabalho necessitava — e que, acreditamos, o tornou substancialmente mais rigoroso
e relevante.

\bigskip\bigskip

\noindent Atenciosamente,\\[0.3em]
\noindent Os Autores

\end{document}
